\dapso{(a) Biogas thải ít \ce{CO2} hơn nếu cung cấp cùng năng lượng; (b) khoảng $70\,200$ đồng/tháng; (c) không bật/tắt thiết bị điện, khóa van gas nếu an toàn, mở cửa thông thoáng và rời khỏi khu vực nếu mùi gas mạnh.}
\begin{loigiaiBT}
Với biogas, $1$ mol \ce{CH4} tỏa \SI{890}{\kilo\joule} và tạo $1$ mol \ce{CO2}. Số mol \ce{CO2} trên mỗi kJ là $1/890$.

Với $1$ mol hỗn hợp LPG:
\[
Q=0.40\times2220+0.60\times2850=\SI{2598}{\kilo\joule},
\]
\[
n_{\ce{CO2}}=0.40\times3+0.60\times4=\SI{3.6}{\mole}.
\]
Do $3.6/2598>1/890$, LPG tạo nhiều \ce{CO2} hơn trên cùng năng lượng.

Khối lượng mol trung bình của LPG:
\[
\overline M=0.40\times44+0.60\times58=\SI{52.4}{\gram\per\mole}.
\]
Số mol LPG trong bình:
\[
n=\frac{12000}{52.4}\approx\SI{229.01}{\mole}.
\]
Tổng nhiệt lượng tỏa ra:
\[
Q=229.01\times2598\approx\SI{594962}{\kilo\joule}.
\]
Nhiệt có ích:
\[
Q_{\text{ích}}=594962\times0.70\approx\SI{416473}{\kilo\joule}.
\]
Số ngày dùng một bình:
\[
N=\frac{416473}{1500}\approx277.65\text{ ngày}.
\]
Chi phí trung bình trong 30 ngày:
\[
C=650000\times\frac{30}{277.65}\approx70200\text{ đồng}.
\]
Khi phát hiện mùi gas, không bật hoặc tắt công tắc điện và không tạo tia lửa; khóa van bình gas nếu có thể thực hiện an toàn, mở cửa để thông thoáng và rời khỏi khu vực nếu mùi gas mạnh.

\textbf{Vậy,} biogas phát thải ít \ce{CO2} hơn trên cùng nhu cầu năng lượng; chi phí gas trung bình khoảng $70\,200$ đồng/tháng; khi rò rỉ gas phải ưu tiên loại bỏ nguồn rò, thông gió và tránh mọi nguồn gây tia lửa.
\end{loigiaiBT}
